\scanpage{32}
% archive foliation: 31
\uline{Corollaire 2} \uline{Dans} $\mathrm{Hot}(\underline W)$, les produits finis existent, et les foncteurs canoniques $(\mathrm{Cat})\longrightarrow\mathrm{Hot}(\underline W)$, $\widehat A\longrightarrow\mathrm{Hot}(\underline W)$ ($A$ catégorie test stricte) y commutent.

\margin{\diagram{\begin{tikzcd}[column sep=small,row sep=small]
{}\arrow[dr]&\\&{}
\end{tikzcd}}il suffit $=$ que $A$ soit tot. $\underline W$-asph.}

On est ramené à prouver que
\[
\widehat A\longrightarrow[\mathrm{Hot}(\underline W)\simeq]\underline W_A^{-1}\widehat A
\]
commute aux produits finis, i.e. que
\[
(*)\quad\operatorname{Hom}_{\underline W_A^{-1}\widehat A}(X,\unsure{Y_j|Y})\longrightarrow\prod_J\operatorname{Hom}_{\underline W_A^{-1}\widehat A}(X,Y_j)
\]
est bijectif, pour \unsure{tte} famille finie $(Y_j)_{j\in J}$ d'objets de $\widehat A$, et $X$ de $\widehat A$.

\margin{\deleted{\unsure{on voit}} Mais le cas de $(\mathrm{Cat})$ se traite directement.}

\uline{Surjectivité} Un élément du $2^{\mathrm o}$ \unsure{mb} provient d'une famille
\diagram{\begin{tikzcd}
f_j:X\arrow[r]&Y'_j
\end{tikzcd}}
où $Y_j\longrightarrow Y'_j$ équiv. faible.

\deleted{\uline{\unsure{Dém}}}\inserted{Considérons} alors
\[
f=(f_j):X\longrightarrow\prod Y'_j=Y'
\]
il définit $X\longrightarrow Y$ qui convient, \uline{si on sait que} $\prod Y_j\longrightarrow\prod Y'_j$ est $\in\underline W$.
